% p146_1; source PDF page 146; manually verified editable table.
\begin{tikzpicture}[x=1cm,y=-1cm,inner sep=0pt,outer sep=0pt]
\path[use as bounding box] (0,0) rectangle (16.000,2.915);
\draw[line width=.4pt] (0.000,0.000) rectangle (2.100,2.915);
\node[align=center,text width=1.840cm,font={\linespread{1}\fontsize{10}{13.50}\selectfont\bfseries}] at (1.050,1.458) {教唆犯的\\成立条件};
\draw[line width=.4pt] (2.100,0.000) rectangle (4.500,2.010);
\node[align=center,text width=2.140cm,font={\linespread{1}\fontsize{10}{13.50}\selectfont}] at (3.300,1.005) {客观要件};
\draw[line width=.4pt] (4.500,0.000) rectangle (8.700,2.010);
\node[align=center,text width=3.940cm,font={\linespread{1}\fontsize{10}{13.50}\selectfont}] at (6.600,1.005) {引起正犯制造违法事实};
\draw[line width=.4pt] (8.700,0.000) rectangle (16.000,1.005);
\node[align=left,text width=7.040cm,font={\linespread{1}\fontsize{10}{13.50}\selectfont}] at (12.350,0.502) {引起正犯制造违法行为，是教唆犯的成立条件};
\draw[line width=.4pt] (8.700,1.005) rectangle (16.000,2.010);
\node[align=left,text width=7.040cm,font={\linespread{1}\fontsize{10}{13.50}\selectfont}] at (12.350,1.507) {引起正犯制造违法结果，是教唆犯的既遂条件};
\draw[line width=.4pt] (2.100,2.010) rectangle (4.500,2.915);
\node[align=center,text width=2.140cm,font={\linespread{1}\fontsize{10}{13.50}\selectfont}] at (3.300,2.462) {主观要件};
\draw[line width=.4pt] (4.500,2.010) rectangle (16.000,2.915);
\node[align=center,text width=11.240cm,font={\linespread{1}\fontsize{10}{13.50}\selectfont}] at (10.250,2.462) {故意引起};
\end{tikzpicture}
